%%%%%%%%%%%%%%%%%%%%%%%%%%%%%%%%%%%%%%%%%%%%%%%%%%%%%%%%%%%%%%
%% Elemento opcional (Guia UECE 2026, 2.2.4.2; NBR 6029).    %%
%% É uma lista, em ordem alfabética, de palavras ou          %%
%% expressões técnicas de uso restrito ou de sentido obscuro %%
%% usadas no trabalho, acompanhadas de seus significados --  %%
%% não é lugar para conceitos já bem conhecidos na sua área. %%
%%                                                            %%
%% Cada entrada usa \newglossaryentry{chave}{name=Termo,      %%
%% description={definição}}. No corpo do texto, use           %%
%% \Gls{chave} (ou \gls{chave}) para remeter a um termo já     %%
%% definido aqui, em vez de repetir a definição no meio do     %%
%% texto.                                                      %%
%%                                                            %%
%% Se o seu trabalho não usar termos técnicos obscuros o      %%
%% suficiente para justificar um glossário, remova esta       %%
%% seção e a linha \imprimirglossario do documento.tex.        %%
%%%%%%%%%%%%%%%%%%%%%%%%%%%%%%%%%%%%%%%%%%%%%%%%%%%%%%%%%%%%%%

\newglossaryentry{exemplo-de-termo}{
	name=Exemplo de Termo,
	description={substitua por um termo técnico do seu trabalho e pela sua definição}
}
